\documentclass[10pt,conference]{IEEEtran}
\usepackage{booktabs}
\usepackage{hyperref}
\usepackage{xurl}

\title{A Replayable Harness for Evaluating Stateful Agent Contracts}
\author{\IEEEauthorblockN{Song Luo}
\IEEEauthorblockA{\texttt{luosongred@gmail.com}}}

\begin{document}
\maketitle

\begin{abstract}
Stateful AI agents increasingly edit repositories, inspect files, call tools,
manage memory, and prepare external side effects. Final-answer evaluation misses
whether such agents used allowed tools, changed the correct state, preserved
replay evidence, or crossed a permission boundary. This demonstration presents a
replayable harness for evaluating production-oriented agent contracts. The
artifact contains 24 pilot cases, mock stateful tools, permission gates, prompt
variants, raw model-run logs, regenerated tables, and validation scripts. The
demo shows how a reviewer can run the harness, inspect state diffs and
permission violations, and regenerate the paper tables from raw JSONL. Tool:
\url{https://github.com/rrrrrredy/agent-harness-paper}. Archive:
\url{https://doi.org/10.5281/zenodo.20907471}. Video: TBD.
\end{abstract}

\section{Motivation}

An agent that produces the right final answer may still fail operationally: it
may read an inappropriate file, call a forbidden tool, ignore untrusted
observations, change the wrong state, or skip human review. These failures are
not visible in a chat-style answer. They are visible only when the evaluation
captures trajectory, tool calls, observations, permissions, and state diffs.

This tool demonstration turns the design principle ``thin harness, strong
contracts'' into an executable artifact. The harness stays thin around model
cognition: it does not impose a heavy planner or role choreography. It is strong
at system boundaries: tool schemas, permission policy, replay logs, memory
scope, and state assertions are explicit and testable.

\section{Artifact}

The repository contains a small but reviewable benchmark and harness:

\begin{itemize}
\item \texttt{benchmark/cases.jsonl}: 24 cases covering routing, stateful tool
use, permission and prompt-injection behavior, memory scoping, and replay.
\item \texttt{harness/}: a mock execution environment, tool simulation, provider
adapter boundary, prompt variants, and metric computation.
\item \texttt{experiments/raw/}: deterministic reference rows and live provider
rows.
\item \texttt{results/tables/}: regenerated summaries, category breakdowns,
failure taxonomy, Wilson intervals, and paired pilot deltas.
\item \texttt{docs/}: reproducibility, claim boundaries, reviewer guide, and
artifact release notes.
\end{itemize}

The artifact is deliberately not a provider leaderboard. It demonstrates an
evaluation shape for stateful agent contracts.

\section{Demo Scenario}

The demonstration follows the path a reviewer can repeat locally.

\begin{enumerate}
\item Validate the artifact:
\texttt{python scripts/run\_checks.py --mode validate}.
\item Run the deterministic reference policy to confirm simulator mechanics.
\item Inspect one security case where untrusted observations try to induce a
forbidden secret read.
\item Regenerate aggregate tables from raw JSONL using the analysis scripts.
\item Compare \texttt{no\_harness}, \texttt{thick\_checklist}, and
\texttt{thin\_contract} prompt variants.
\item Inspect the GitHub and Zenodo release state.
\end{enumerate}

The video highlights the reviewer workflow rather than only showing slides:
case definition, command execution, generated metrics, and a concrete
state-diff/permission failure.

\section{Use Cases}

The harness is useful when an agent workflow needs evidence that work was
completed by changing the correct state under the correct policy. Examples
include coding agents that should modify a repository and pass tests, research
agents that should preserve citation provenance, memory agents that must avoid
cross-project recall, and workflow agents that should dry-run or escalate
irreversible actions.

\section{Pilot Evidence}

The pilot contains 24 cases and three prompt variants. A deterministic reference
runner checks case and metric consistency. Two live API runs provide 72 rows per
provider. In the committed pilot rows, thin-contract prompting improves task
success over a thick checklist for both providers, but these numbers are
descriptive observations about this artifact, not general claims about model
quality. The important output is the failure structure: invalid tool calls,
permission violations, parser failures, missing state diffs, and recovery
errors become inspectable.

\section{Availability}

The tool and data are publicly available at
\url{https://github.com/rrrrrredy/agent-harness-paper}. The archival artifact is
available at \url{https://doi.org/10.5281/zenodo.20907471}. The repository is
licensed under MIT for code and CC BY 4.0 for manuscript, documentation,
benchmark cases, evidence snapshots, tables, and figures.

\begin{thebibliography}{9}
\bibitem{toolsandbox} ToolSandbox, stateful tool-use evaluation benchmark.
\bibitem{agentdojo} AgentDojo, dynamic tool-calling and prompt-injection evaluation.
\bibitem{sweagent} SWE-agent and agent-computer interfaces for software engineering agents.
\bibitem{deli} Deli AutoResearch framework for long-horizon autonomous tasks.
\end{thebibliography}

\end{document}

